\chapter{计算几何}

\section{基础几何}

\subsection{点、直线与常用运算}

用 \texttt{Point<T>\{x,y\}} 和 \texttt{Line<T>\{a,b\}} 构造点与直线。

用 \texttt{dot(a,b)}、\texttt{cross(a,b)} 和 \texttt{dis(a,b)} 求点积、叉积与距离。需要分数坐标时，点类型使用浮点数。

\cppfile{code/geometry/g1.cpp}

\subsection{曼哈顿距离与切比雪夫距离}

求非空点集的最大曼哈顿距离时，分别计算 \texttt{x+y} 与 \texttt{x-y} 的最大值减最小值，再取较大者。

\begin{pseudocode}
u = 所有点的 x+y
v = 所有点的 x-y
ans = max(max(u)-min(u), max(v)-min(v))
\end{pseudocode}

\section{多边形}

\subsection{多边形与圆相交面积}

给定半径 \texttt{R}，将多边形顶点平移到以圆心为原点的坐标系，并沿边界顺序放入 \texttt{q}。调用 \texttt{work(q)} 返回交集面积。

\cppfile{code/geometry/g2.cpp}

\subsection{二维凸包}

对非空点集调用 \texttt{getConvex(a)}。它会重排 \texttt{a}，返回排序后数组中的凸包下标，末尾包含重复的起点。

\cppfile{code/geometry/g3.cpp}

\subsection{共线情形的凸包周长}

准备 \texttt{getConvex(a)} 返回的 \texttt{res}，以及点数 \texttt{n}。执行下列片段，周长存于 \texttt{ans}。

\cppfile{code/geometry/g3-1.cpp}

\subsection{三维凸包}

输入点数 \texttt{n}，随后逐行输入三维点坐标。至少给出四个非共面的点。程序输出凸包表面积，保留六位小数。

\cppfile[firstline=4]{code/geometry/g4.cpp}

\section{常考进阶}

\subsection{最小覆盖圆}

将非空点集放入 \texttt{q}，调用 \texttt{MinimumCircleCovering(q)} 返回包含所有点的圆。函数会调整 \texttt{q} 中的点序。

\cppfile{code/geometry/g5.cpp}

\subsection{最小覆盖矩形}

取非退化凸包的闭合下标序列 \texttt{ids} 和对应的浮点坐标数组 \texttt{q}。

\texttt{MinimumRectangularCover(ids,q)} 返回最小面积和矩形四个顶点。

\cppfile{code/geometry/g6.cpp}

\subsection{旋转卡壳}

用 \texttt{getConvex(a)} 得到闭合凸包下标 \texttt{id}。调用 \texttt{RotatingCalipers(id,a)} 返回最远点对距离，\texttt{a} 使用求凸包排序后的数组。

\cppfile{code/geometry/g7.cpp}

\subsection{半平面交}

用 \texttt{Line<T>} 存半平面的有向边界，取每条边的左侧为有效区域。调用 \texttt{hp(lines)}，返回有界交集的顶点。只需要在给定范围内求交时加入外包矩形的四条边。

\cppfile{code/geometry/g8.cpp}
